\section{An actual quartic-character ordinary seed at the original coefficient prime}
\label{sec:cubic-quartic-actual-seed}

This section constructs a nonempty arithmetic input after the quadratic-character
exclusion.  The Dirichlet modulus in the source datum, the level of the initial
Eisenstein-congruent newform, and the level of the endpoint newform are different:
\[
                 N_\psi=10,\qquad N_g=15,\qquad N_{f_0}=75.
                                                        \tag{G118.1}
\]
No class-sector estimate from the earlier quadratic example is transferred to
this datum.  The local selected-component assertion retains the source
\texttt{LocalPublishedTheorems} input, as in AC-G48-TypeII.

Put $p=269$, $K=\mathbb Q(\sqrt{-3})$, and let $c_p$ denote the actual
action on $\mu_{p^\infty}$, so that $c_p(\operatorname{Fr}_{q,\mathrm{arith}})=q$.
Write $\eta=\chi_{-3}$, both for the Dirichlet character and its Galois avatar
when the domain is clear.  Choose $i\in\mathbb Z_p$ with $i^2=-1$ and
$i\equiv82\pmod p$.  Let $\psi$ be the primitive quartic character of conductor
$5$ with $\psi(2)=i$, viewed as a character of modulus $10$.  Its Galois avatar
has arithmetic-Frobenius value $\psi(q)$ at $q\ne5$, where this value and every
automorphic twisting character below use the primitive conductor-$5$ avatar.
The modulus $10$ is used only for the source datum and has the same Galois
character; in particular no imprimitive zero at $2$ is used.  Set
\[
 e=23,\quad k=2e+1=47,\quad \mathcal O_0=\mathbb Z_p,\quad
 \theta=\eta\psi^{-2}=\chi_{-15},\qquad \Psi=\psi c_p^{23}.
                                                        \tag{G118.2}
\]
In particular $\theta$ is primitive of conductor $15$, is odd, and takes
values in $\{\pm1\}$.

\begin{theorem}[AC-G118-quartic-seed: an actual original ordinary Type II seed]
\label{AC-G118-quartic-seed}
With (G118.1)--(G118.2), the following objects exist.

\begin{enumerate}
\item There is a normalized cuspidal newform $g$ of weight $47$, primitive
level $15$, and nebentype $\theta$, together with a place over $269$, such that
\[
 a_q(g)\equiv1+\theta(q)q^{46}\pmod{\mathfrak m}
       \quad(q\nmid15),\qquad a_3(g)\equiv a_5(g)\equiv1\pmod{\mathfrak m}.
                                                        \tag{G118.3}
\]
In particular $a_p(g)$ is a unit.

\item Let $R_g$ be its cohomological Galois representation: its
characteristic polynomial on geometric Frobenius at a good prime $q$ is
$X^2-a_q(g)X+\theta(q)q^{46}$.  Then
\[
       V=R_g\otimes\eta\psi c_p^{23},\qquad \det V=\eta,
                                                        \tag{G118.4}
\]
is absolutely irreducible on $G_{\mathbb Q}$ and on $G_K$.
It is the normalized representation of the primitive automorphic twist
\[
                    f_0=g\otimes(\eta\psi^{-1}),
                                                        \tag{G118.5}
\]
which has weight $47$, nebentype $\eta$, primitive level $75$, and is ordinary
at $269$.  The inverse on $\psi$ in (G118.5) is required by the geometric
Frobenius convention.  After a finite extension of the value field, $V$ has
a stable lattice whose residual representation is the split representation
\[
             \overline V=
             (\eta\overline\psi\,\overline c_p^{23})
                \oplus(\overline\psi^{-1}\overline c_p^{-23}).
                                                        \tag{G118.6}
\]
The original coefficient ring $\mathcal O_0$ is unchanged.

\item For the actual old datum $D=\texttt{Dat}$ associated with
$(K,p,e,\psi,\mathcal O_0)$ and a marked pair $(c,g_1)$ as constructed below,
there is a continuous integral homomorphism
\[
                       r:G_{\mathbb Q}\longrightarrow GG_3(\overline{\mathbb Q}_p)
                                                        \tag{G118.7}
\]
with constant middle line $1$, endpoint representation $V$ on $G_K$, and
the original split residual representation $r_0(\eta,\Psi)$.
It is \texttt{TypeII} and satisfies all seven fields of
$D.\texttt{IsPt}(U,g_1,r)$ for an actual framed local universal object $U$.
In particular its local point kills the original selected ideal $I_t$.
For every actual source global universal object $S:D.\texttt{Ring}(U,g_1)$,
it gives an $\operatorname{IsLocPt}$ specialization of $S$.
This is membership in the original $I_t$; no new equality with a cubic
selected prime $P^*$ is asserted.
\end{enumerate}

The automorphic and local compatibility theorems used in the proof are
retained published inputs.  No prior old point, no auxiliary-level seed,
no host class-group estimate, and no mixed Selmer vanishing are assumed.
\end{theorem}

\begin{proof}
\proofstep{1. The exact constant and every cusp.}
Use the Eisenstein normalization
\[
 E=E_{47}(1,\theta)
  =c+\sum_{n\ge1}\left(\sum_{d\mid n}\theta(d)d^{46}\right)q^n,
 \qquad c=\frac{L(-46,\theta)}2=-\frac{B_{47,\theta}}{94}.
                                                        \tag{G118.8}
\]
The following is an exact integer calculation (the full integers are also
in the machine-readable receipt):
{\small\begin{align}
 B_{47,\theta}
 &=77176891038967163668552282889941348327824135882255355566505306467596555167952,\notag\\
 c
 &=-821030755733693230516513647765333492849192934917610165601120281570176118808,\notag\\
 v_{269}(c)&=1,\qquad (c/269)\bmod269=188. \tag{G118.9}
\end{align}}
For reproducibility, use $B_0=1$ and
\[
 B_m=-\frac1{m+1}\sum_{a=0}^{m-1}\binom{m+1}{a}B_a,
 \qquad
 B_{47,\theta}=15^{46}\sum_{a=1}^{15}\theta(a)B_{47}(a/15).
                                                        \tag{G118.10}
\]
All operations in the saved script are rational operations before reduction.
An independent rational Akiyama--Tanigawa calculation gives the same integer
and valuation.  This calculation alone is not used as an existence theorem.

The all-cusp assertion follows from Dasgupta--Kakde,
\emph{On constant terms of Eisenstein series}, Definitions 4.1--4.4 and
Theorem 4.5(1), with its normalization in section 4.  Apply it with
$\chi_1=1$, $\chi_2=\theta$, $\mathfrak b=\mathfrak n=\mathfrak f=(15)$.
The Gauss-sum quotient in (30) is $1$, the norm quotient is $1$, and
$T_{\mathfrak n,\mathfrak f}$ is empty.  For an integral determinant-one
cusp representative with primitive first column $(a,c_0)$, the constant is
zero unless $15\mid c_0$.  In the latter case it is a sign times
$\theta(a)^{-1}c$, hence a $269$-adic unit times $c$.
The convention for $c_0=0$ gives the same assertion at infinity.
Thus every cusp constant is divisible by $269$.
Pullback to full level $15$ changes the cusp parameter but does not change
the vanishing of its constant; all level-normalizing factors are $269$-adic
units.  This uses the primitive nontrivial character $\chi_{-15}$ itself,
not a trivial-character constant-term formula.

\proofstep{2. Integral cusp lifting in the exact character space.}
Let $\mathcal O$ be the unramified extension of $\mathbb Z_p$ containing
$\zeta_{15}$, and choose a determinant component $X$ of the compact fine
full-level-$15$ modular curve over $\mathcal O$.
Here $[\mathcal O:\mathbb Z_p]=2$ because $269\equiv-1\pmod {15}$.
Let $C$ be its reduced cusp divisor and $\omega$ its invariant-differential
line bundle.  This curve is smooth and proper at $269$; the cusp divisor is
finite etale there.  These are the tame integral modular-curve models of
Katz, \emph{$p$-adic properties of modular schemes and modular forms},
sections 1.5--1.7.

Katz's q-expansion principle, Corollary 1.6.2, applied on this determinant
component at its standard Tate cusp, puts $E$ in
$H^0(X,\omega^{47})$.  Its q-coefficients and its constant in (G118.9) are
integral.  The preceding all-cusp calculation gives
\[
       0\ne\overline E\in
       H^0(X_{\mathcal O/p},\omega^{47}(-C)),
                                                        \tag{G118.11}
\]
where nonzero follows from the coefficient of $q$, which is $1$.

For completeness the needed cusp-form base change follows from the same
Kodaira--Spencer identity used in Katz's full proof of Theorem 1.7.1:
\[
                 \omega^2\simeq\Omega^1_X(C).
                                                        \tag{G118.12}
\]
On each geometric special-fiber component, Serre duality identifies the
dual of $H^1(\omega^{47}(-C))$ with $H^0(\omega^{-45})$.
The degree of $\omega$ is positive (equivalently the discriminant section
of $\omega^{12}$ has its nonempty effective cusp divisor).
Hence the latter space is zero.  The exact reduction sequence and
cohomology-and-base-change give a surjection
\[
 H^0(X,\omega^{47}(-C))\longrightarrow
 H^0(X_{\mathcal O/p},\omega^{47}(-C)).
                                                        \tag{G118.13}
\]
The source is finite and torsionfree over the DVR $\mathcal O$, hence free.

The group imposing the classical $\Gamma_0(15)$ condition on this full-level
component is the upper Borel
\[
 \mathcal B=\left\{\begin{pmatrix}a&b\\0&d\end{pmatrix}
       \in\mathrm{SL}_2(\mathbb Z/15\mathbb Z)\right\},
       \qquad |\mathcal B|=15\varphi(15)=120.
                                                        \tag{G118.14}
\]
Use the usual slash action on sections, and the character
$\chi_{\mathcal B}(\begin{smallmatrix}a&b\\0&d\end{smallmatrix})=\theta(d)$.
The projector
\[
 e_\theta=\frac1{120}\sum_{b\in\mathcal B}\chi_{\mathcal B}(b)^{-1}b
                                                        \tag{G118.15}
\]
is integral, preserves cusp forms, and commutes with reduction.
The central element $-I$ acts on the weight line by $(-1)^{47}$, agreeing
with $\theta(-1)$; no freeness of the Borel action is assumed.
Consequently (G118.13) is surjective on the actual character eigenspace
$S_{47}(\Gamma_0(15),\theta;\mathcal O)$, and (G118.11) belongs to its
reduction.  This uses only a denominator $120$, not an order-$269$
character projector.

\proofstep{3. Eigenvalues lift to an actual primitive ordinary newform.}
The commuting integral Hecke operators preserve this finite free cusp module.
The integral q-expansion formulas show that $\overline E$ is a simultaneous
eigenvector, with eigenvalues
\[
 T_q:1+\theta(q)q^{46}\ (q\nmid15),\qquad U_3:1,\quad U_5:1.
                                                        \tag{G118.16}
\]
For $q=3,5$, the equality follows directly from the divisor sum in (G118.8),
since $\theta(q)=0$; the constant also has the required $U_q$ action.
Apply the full statement of Deligne--Serre, Lemma 6.11, to these commuting
operators.  After a finite extension of the value field there is a nonzero
characteristic-zero simultaneous cusp eigenvector with these reduced
eigenvalues.  This lemma lifts the eigenvalues, not necessarily the given
vector $\overline E$.  The family of operators may be infinite: the generated
algebra is a finite $\mathcal O$-submodule of the endomorphism algebra.

Classical newform decomposition gives a normalized primitive newform in this
eigensystem with primitive level dividing $15$.  Its central character is
the primitive character $\theta$ of conductor $15$, so its level is exactly
$15$; no old subspace from a proper divisor can have that character.
It has (G118.3), including the bad-prime eigenvalues.
In particular $a_p(g)\equiv1$, so $g$ is ordinary.
Only the existence of its Hecke eigensystem is asserted here, not a numerical
enumeration of all its characteristic-zero coefficients.

\proofstep{4. Galois and primitive-twist normalizations.}
Use the published attachment theorem for classical cuspidal newforms,
including absolute irreducibility, crystallinity at good coefficient primes,
and local-global Weil--Deligne compatibility with monodromy.
In the source convention $R_g$ is the dual of the representation with
arithmetic-Frobenius Hecke polynomial.  Therefore
\[
 \det R_g=\theta^{-1}c_p^{-46}=\theta c_p^{-46},\qquad
 \overline R_g^{\mathrm{ss}}=1\oplus\theta\overline c_p^{-46}.
                                                        \tag{G118.17}
\]
For example the trace of this right-hand side on arithmetic Frobenius is
$1+\theta(q)q^{-46}$, obtained by dualizing
$1+\theta(q)q^{46}$ with determinant $\theta(q)q^{46}$.
Chebotarev and Brauer--Nesbitt give (G118.17) from (G118.3).

At $3$ and $5$, the primitive local conductor exponent is one and the central
character has conductor exponent one.  The local parameter is a sum of one
unramified character and the prescribed conductor-one character, and $N=0$.
Indeed a conductor-one special representation has unramified central character;
a supercuspidal parameter has conductor at least two.  Thus the inertia of
$R_g$ is $1\oplus\eta$ at $3$, and $1\oplus\psi^{-2}$ at $5$.
It is unramified at every other prime different from $269$.

At $p$ it is crystalline.  Its crystalline Frobenius polynomial has roots
of valuations $0,46$, since $a_p(g)$ is a unit and its determinant has
valuation $46$.  Its filtration jumps are $0,46$.
The unique slope-zero line has Hodge degree at least zero, while weak
admissibility bounds it above by its Newton degree zero.  It is therefore
weakly admissible, and the crystalline equivalence gives an unramified
Galois subline.  Its quotient has inertia $c_p^{-46}$.
This also proves the direction of the ordinary flag used below.

Define $V$ by (G118.4).  Equations (G118.2) and (G118.17) give its determinant
and residual semisimplification (G118.6).  Its ordered inertia characters
at $p$ are $\psi c_p^{23}$ and $\psi^{-1}c_p^{-23}$, with $\psi$ unramified
there.  Away from $p$ its finite inertia is
\[
 \begin{array}{c|c|c}
 \text{place}&V|_I^{\mathrm{ss}}&N\\ \hline
 3&1\oplus\eta&0\\
 5&\psi\oplus\psi^{-1}&0\\
 q\ne3,5,p&1\oplus1&0.
 \end{array}                                             \tag{G118.18}
\]
These are actual finite inertia representations, not assertions deduced just
from their residual semisimplifications.

At geometric Frobenius the factor $\eta\psi$ in (G118.4) is
$\eta(q)\psi(q)^{-1}$.  Hence $V c_p^{-23}$ is attached in the literal
\texttt{Newform.Att} convention to (G118.5).
That primitive automorphic twist has nebentype
$\theta(\eta\psi^{-1})^2=\eta$.
At $3$ twisting interchanges the unramified and ramified characters, leaving
conductor exponent one.  At $5$ both resulting characters have conductor
exponent one, giving exponent two; $N=0$ is preserved by the character twist.
No other prime is introduced.
Thus its primitive level is exactly $75$.
At $p$, $\eta(p)\psi(p)^{-1}=1$, so its ordinary eigenvalue is even the same
as that of $g$.  In particular it is a good-coefficient-prime ordinary form.

\proofstep{5. Absolute irreducibility after restriction to $K$.}
Absolute irreducibility on $G_{\mathbb Q}$ is the classical cuspidal attachment
theorem and is preserved by scalar twisting.  If its restriction to the normal
index-two subgroup $G_K$ were reducible, Clifford theory and Frobenius
reciprocity would imply $V\simeq V\otimes\eta$.
For clarity, restriction is semisimple: translates of a simple $G_K$-submodule
sum to a nonzero $G_{\mathbb Q}$-submodule and hence to the whole space.
The identity
$\operatorname{End}_{G_K}(V)=
 \operatorname{Hom}_{G_{\mathbb Q}}(V,V\oplus V\eta)$
then yields this self-twist.

It would also give $R_g\simeq R_g\eta$.
The unique inertia-trivial subline of the ordinary $R_g|_{G_{\mathbb Q_p}}$
must map to itself: its other inertia character is $c_p^{-46}\ne1$.
If the unramified subcharacter is $u$, this would give $u=u\eta_p$.
But $269$ is inert in $K$, so $\eta_p$ is the nontrivial unramified quadratic
character.  This is impossible.  Thus $V|_{G_K}$ is absolutely irreducible.

\proofstep{6. An actual split residual lattice without replacing $\mathcal O_0$.}
Fix a finite value field containing all coefficients of $V$, and a stable lattice.
Its residual semisimplification has the two distinct characters in (G118.6),
defined over $\mathbb F_p$.  Choose a basis making its residual matrices upper
triangular.  Write the matrices as $\left(\begin{smallmatrix}a&b\\c&d\end{smallmatrix}\right)$
with $b\in\mathcal O_E$, $c\in\varpi_E\mathcal O_E$.
Pass to a ramified quadratic extension $E'/E$ with
$\varpi_E$ a unit times $\pi'^2$ and conjugate by
$\operatorname{diag}(\pi',1)$.
The off-diagonal entries become $\pi'b$ and $c/\pi'$, both in
$\pi'\mathcal O_{E'}$.  The resulting lattice is stable and has split residual
representation.  Reorder its two lines if necessary to give (G118.6).
This procedure enlarges only the value field; it does not assert a split
residual lattice over the original value DVR.

Complex conjugation has square one and reduces to $D_2=\operatorname{diag}(-1,1)$.
The columns $(1-V(c))/2$ applied to the first basis vector and
$(1+V(c))/2$ applied to the second give a basis congruent to the existing one
in which $V(c)=D_2$ exactly.
Every matrix entry now has residue in $\mathbb F_p$.  Set
\[
                  A=\mathcal O_0+\mathfrak m_{E'}\subset\mathcal O_{E'}.
                                                        \tag{G118.19}
\]
This finite complete local $\mathcal O_0$-algebra has maximal ideal
$\mathfrak m_{E'}$ and residue field $\mathbb F_p$, and contains all matrices
and their inverses.  Its maximal-ideal topology agrees with the restricted
valuation topology: if $a=v_{E'}(p)$, then
$\mathfrak m_A^{a+1}\subset pA\subset\mathfrak m_A$.
This is the same original-coefficient construction used in AC-K29-point.

\proofstep{7. Literal datum, multiplier, normalization, and away killing.}
The source \texttt{GD} clauses hold: $p$ is odd, $e>0$, $N_\psi=10$ is even,
$\operatorname{ord}\psi=4$ is prime to $p$, $\operatorname{comp}_p\psi=1$,
$\psi(-1)=-1=(-1)^{23}$, $K$ is imaginary quadratic,
$p\nmid3h_K$ with $h_K=1$, $\gcd(3,10)=1$, and $268\nmid46$.
The remaining \texttt{GD.eta} clause follows because $K$ is ramified at $3$,
whereas the field cut out by $\overline c_p$ and $\psi$ is unramified at $3$.
The local number field in \texttt{Dat.Kv} is the unramified quadratic
$K_p/\mathbb Q_p$, with $\eta$ trivial on coefficient-prime inertia.
All character-continuity and \texttt{Psi\_I} clauses follow from the actual
cyclotomic character and the unramified local character $\psi$.

Choose a complex conjugation $c$ and a local inertia element $t$ with
$c_p(t)=2$, and let $g_1$ be its global image through the chosen local embedding.
Then $\eta(g_1)=1$, $\Psi(c)=1$, and
$\overline\Psi(g_1)^2=2^{46}\ne1\pmod{269}$.
The last inequality follows for example from $2^{134}=-1$ and $2^4\ne1$
modulo $269$, so $2$ has order $268$.
Thus every literal field of \texttt{Marks} holds.

Write $\eta(g)=(-1)^{\epsilon(g)}$ and put $B(g)=V(g)D_2^{\epsilon(g)}$.
Then $\det B(g)=1$.  Embed $B(g)$ in coordinates $0,2$ and put $1$ in the
middle coordinate.  This is the matrix component of $r(g)$, with multiplier
$\eta(g)$.  The identity
\[
             J_2B^{-t}J_2=D_2BD_2\qquad(B\in\mathrm{SL}_2)
                                                        \tag{G118.20}
\]
proves the semidirect multiplication law.  It is an $A$-valued continuous
homomorphism, integral, and its reduction is exactly $r_0(\eta,\Psi)$.
Since $V(c)=D_2$, its matrix at $c$ is $I$.
All mixed entries vanish; hence the two mixed entries of
$r(g_1)-r(cg_1^{-1}c)$ vanish as well.  These are precisely the three
\texttt{Normalised} fields, without an additional unproved diagonalization.

If $g\in D.\texttt{Kill}$, it is an away-$p$ inertia element with residual
multiplier one and residual $\psi(g)=1$.
Reduction is faithful on the order-four image of $\psi$, so $\eta(g)=\psi(g)=1$.
Equation (G118.18) now gives $V(g)=I$, at $3$ as well as at $5$ and the
unramified primes.  Thus $r(g)=1$.
The finite inertia and $N=0$ established above are needed here; residual
unramification alone would not prove this field.

\proofstep{8. The actual local selected ideal and the source ring map.}
At $p$, the middle line is constant and the endpoint block has determinant
one everywhere.  It is crystalline on
$H_L=\ker\eta\cap\ker(\operatorname{comp}_p\psi)=G_{K_p}$.
The stable ordinary line of Step 4, followed by the middle line, gives an
actual ordered complete flag on $\ker\eta$, with inertia characters
\[
                          \Psi,\quad1,\quad\Psi^{-1}.
                                                        \tag{G118.21}
\]
Their endpoint ratio has residual value $\overline c_p^{46}\ne1$ on inertia.
AC-G48-TypeII therefore applies and proves that the local point kills the
original selected ideal $I_t$ over the original $\mathcal O_0$.
Its finite coefficient-descent argument is part of that result; there is no
assumption that selected primes commute with arbitrary base change.

As throughout this source interface, $U$ is the actual unrestricted framed
local universal object; its existence uses local framed representability
(continuous $H^1(F,\mathbb F_p)$ is finite for every finite extension
$F/\mathbb Q_p$).  No new global modularity hypothesis supplies this object.
More explicitly the $A$-valued local framed lift gives, by the actual framed
universal property, a local map $f_U:U.R\to A\subset E'$, and
\[
 \operatorname{IsLocPt}(f_U),\qquad I_t\subseteq\ker f_U,
        \qquad\operatorname{push}(f_U,U.r)=r\circ\iota.
                                                        \tag{G118.22}
\]
The exact seven fields of \texttt{Dat.IsPt} in
\path{UnitaryThreeGroups/RhoKApply.lean} are now verified: \texttt{mult},
\texttt{cont}, \texttt{int}, \texttt{red}, \texttt{norm}, \texttt{kill},
and the local factorization \texttt{loc} in (G118.22).
The definition of \texttt{TypeII} is \texttt{Block} together with irreducibility
of the endpoint corner on $\ker\eta$.  The block property is explicit in the
construction, and Step 5 proves the required irreducibility.

Finally $A$ satisfies the literal \texttt{IsCO} record with original residue
field.  The representation satisfies \texttt{NormCond} and \texttt{BoxCond}:
the latter is exactly away killing and the local factorization through $I_t$.
The \texttt{GRepresents} universal property therefore supplies a map
$S.R\to A\subset E'$ for every actual source object $S$.
It is local (the source and target have the same residue field with a
surjective coefficient residue map), hence is the claimed
$\operatorname{IsLocPt}$ specialization.  This completes all point assertions.
\end{proof}

\paragraph{Arithmetic meaning and retained limits.}
The form $f_0$ supplies the actual ordinary modular endpoint required by the
displayed starting hypotheses of AC-G78, with $e=23$ and $\tau=1$; its local
sign is the automatic odd-$e$ sign.  Subsequent host and unitary constructions
may be applied only with their own stated hypotheses.
In particular the existing quadratic-character class certificate is not a
certificate for $\psi$ here.  This theorem neither proves the cubic mixed
Selmer vanishing nor closes a numerical frame inequality, a whole-curve
Hecke-ideal containment, or H04.4.

\paragraph{Source and verification scope.}
The exact Bernoulli computation and its independent rational replay are
saved beside the source receipt.  The complete applied Dasgupta--Kakde
constant-term proof and Deligne--Serre Lemma 6.11 proof were read, as were
Katz's complete sections 1.6--1.7 and the containing tame-curve definitions.
The cusp base-change, Borel projector, lattice splitting, no-self-twist,
semidirect construction, and seven-field application above are proved here.
The attachment theorem is used as a published input, with an exact convenient
restatement in Atkin--Li--Liu--Long, \emph{Galois representations with quaternion
multiplication}, Theorem 2.1.1, printed p.6220.
Saito's author survey, sections 2a and 2c, Theorem 3, supplies the stated full
Weil--Deligne compatibility; Casselman's newform theorem and its global
application supply the primitive conductor interface.  These deep theorems,
the crystalline equivalence, and the source
\texttt{LocalPublishedTheorems} are retained inputs, not re-proved here.
The complete native \texttt{Counts}, \texttt{RhoKApply}, \texttt{Newform},
\texttt{UnitaryThreeReflectionRight}, \texttt{BaseRepGL3NormalisedOfPsiEqOne},
\texttt{GRepresents}, and \texttt{ResMapBijective} files and the applied
AC-G48-TypeII proof are the source interface being checked.
No Lean execution or original modular-form enumeration is claimed.
